%-----------------------------------------------------------------------------------------------------------------------------------------------%
%	The MIT License (MIT)
%
%	Copyright (c) 2021 Jitin Nair
%
%	Permission is hereby granted, free of charge, to any person obtaining a copy
%	of this software and associated documentation files (the "Software"), to deal
%	in the Software without restriction, including without limitation the rights
%	to use, copy, modify, merge, publish, distribute, sublicense, and/or sell
%	copies of the Software, and to permit persons to whom the Software is
%	furnished to do so, subject to the following conditions:
%
%	THE SOFTWARE IS PROVIDED "AS IS", WITHOUT WARRANTY OF ANY KIND, EXPRESS OR
%	IMPLIED, INCLUDING BUT NOT LIMITED TO THE WARRANTIES OF MERCHANTABILITY,
%	FITNESS FOR A PARTICULAR PURPOSE AND NONINFRINGEMENT.
%	IN NO EVENT SHALL THE
%	AUTHORS OR COPYRIGHT HOLDERS BE LIABLE FOR ANY CLAIM, DAMAGES OR OTHER
%	LIABILITY, WHETHER IN AN ACTION OF CONTRACT, TORT OR OTHERWISE, ARISING FROM,
%	OUT OF OR IN CONNECTION WITH THE SOFTWARE OR THE USE OR OTHER DEALINGS IN
%	THE SOFTWARE.
%
%
%-----------------------------------------------------------------------------------------------------------------------------------------------%

%----------------------------------------------------------------------------------------
%	DOCUMENT DEFINITION
%----------------------------------------------------------------------------------------

% article class because we want to fully customize the page and not use a cv template
\documentclass[a4paper,12pt]{article}

%----------------------------------------------------------------------------------------
%	CHINESE SUPPORT & FONT
%----------------------------------------------------------------------------------------

% Added ctex for Chinese support. IMPORTANT: Compile with XeLaTeX.
\usepackage[UTF8]{ctex}

%----------------------------------------------------------------------------------------
%	PACKAGES
%----------------------------------------------------------------------------------------
\usepackage{url}
\usepackage{parskip}

%other packages for formatting
\RequirePackage{color}
\RequirePackage{graphicx}
\usepackage[usenames,dvipsnames]{xcolor}
\usepackage[scale=0.9]{geometry}

%tabularx environment
\usepackage{tabularx}

%for lists within experience section
\usepackage{enumitem}

% centered version of 'X' col.
\newcolumntype{C}{>{\centering\arraybackslash}X}

%to prevent spillover of tabular into next pages
\usepackage{supertabular}
\usepackage{tabularx}
\newlength{\fullcollw}
\setlength{\fullcollw}{0.47\textwidth}

%custom \section
\usepackage{titlesec}
\usepackage{multicol}
\usepackage{multirow}

%CV Sections inspired by:
%http://stefano.italians.nl/archives/26
\definecolor{sectioncolor}{RGB}{0, 102, 204}  % Define blue color for sections
\titleformat{\section}{\Large\scshape\raggedright\color{sectioncolor}}{}{0em}{}[\titlerule]
\titlespacing{\section}{0pt}{10pt}{10pt}

%for publications
\usepackage[style=authoryear,sorting=ynt, maxbibnames=2]{biblatex}

%Setup hyperref package, and colours for links
\usepackage[unicode, draft=false]{hyperref}
\definecolor{linkcolour}{rgb}{0,0.2,0.6}
\hypersetup{colorlinks,breaklinks,urlcolor=linkcolour,linkcolor=linkcolour}
\addbibresource{citations.bib}
\setlength\bibitemsep{1em}

%for social icons
\usepackage{fontawesome5}

%debug page outer frames
%\usepackage{showframe}


% job listing environments
\newenvironment{jobshort}[2]
    {
    \begin{tabularx}{\linewidth}{@{}l X r@{}}
    \textbf{#1} & \hfill &  #2 \\[3.75pt]
    \end{tabularx}
    }
    {
    }

\newenvironment{joblong}[2]
    {
    \begin{tabularx}{\linewidth}{@{}l X r@{}}
    \textbf{#1} & \hfill &  #2 \\[3.75pt]
    \end{tabularx}
    \begin{minipage}[t]{\linewidth}
    \begin{itemize}[nosep,after=\strut, leftmargin=1em, itemsep=3pt,label=--]
    }
    {
    \end{itemize}
    \end{minipage}
    }



%----------------------------------------------------------------------------------------
%	BEGIN DOCUMENT
%----------------------------------------------------------------------------------------
\begin{document}

% non-numbered pages
\pagestyle{empty}

%----------------------------------------------------------------------------------------
%	TITLE
%----------------------------------------------------------------------------------------

\begin{tabularx}{\linewidth}{@{} X r @{}}
\Huge{高恒斌} & \multirow{7}{*}{\includegraphics[width=3cm]{profile.jpg}} \\[3pt]
\normalsize{\textit{本科生}} & \\
\normalsize{\textit{浙江大学—信息工程}} & \\[10pt]
\href{mailto:Particle_Nebula@outlook.com}{\raisebox{-0.05\height}\faEnvelope \ Particle\_Nebula@outlook.com} \ $|$ \
\href{tel:+8619883109386}{\raisebox{-0.05\height}\faMobile \ (+86)\ 19883109386} & \\[3pt]
\href{mailto:3230105132@zju.edu.cn}{\raisebox{-0.05\height}\faEnvelope \ 3230105132@zju.edu.cn} \ $|$ \
\href{https://Particlela.github.io}{\raisebox{-0.05\height}\faGlobe \ Particlela.github.io} & \\[3pt]
\href{https://github.com/Particlela}{\raisebox{-0.05\height}\faGithub\ Particlela} & \\
\end{tabularx}
\vspace{10pt}

%----------------------------------------------------------------------------------------
% EXPERIENCE SECTIONS
%----------------------------------------------------------------------------------------

\section{个人简介}
我是浙江大学信息科学与电子工程学院信息工程专业的大三本科生。作为 RoboMaster 电控组成员，我的科研聚焦于串联弹性腿机器人的控制，并研究了四开关 Buck-Boost 缓冲电容器的控制方法。未来，我计划进一步探索机器学习、数字信号/图像处理与嵌入式系统等方向。欢迎联系我交流！

%----------------------------------------------------------------------------------------
%	EDUCATION
%----------------------------------------------------------------------------------------
\section{教育背景}
\begin{tabularx}{\linewidth}{@{}l X@{}}
2023 - 2027 & 工学学士，\textbf{浙江大学} 信息工程 \hfill \normalsize \\
& 信息科学与电子工程学院 \\
& 专业 GPA：4.48/5.0 \\
& 总 GPA：4.46/5.0 \\
\end{tabularx}

%----------------------------------------------------------------------------------------
%	HONORS AND AWARDS
%----------------------------------------------------------------------------------------
\section{荣誉奖项}
\begin{tabularx}{\linewidth}{@{}l X@{}}
2026 & \textbf{三等奖}，全国大学生机器人大赛 RoboMaster 机甲大师超级对抗赛—步兵机器人竞赛 \\[3.75pt]
2026 & \textbf{二等奖}，全国大学生机器人大赛 RoboMaster 机甲大师超级对抗赛—哨兵机器人竞赛 \\[3.75pt]
2025 & \textbf{一等奖}，全国大学生机器人大赛 RoboMaster 机甲大师超级对抗赛—全国总决赛 \textbf{(全国八强)} \\[3.75pt]
2025 & \textbf{一等奖}，全国大学生机器人大赛 RoboMaster 机甲大师超级对抗赛—步兵机器人竞赛 \\[3.75pt]
2025 & \textbf{三等奖学金}，国家人才培养基地 \\[3.75pt]
2025 & \textbf{三等奖学金}，浙江大学 \\[3.75pt]
2024 & \textbf{浙江省政府奖学金}，浙江大学 \\[3.75pt]
2024 & \textbf{二等奖学金}，浙江大学 \\
\end{tabularx}

%Projects
\section{项目经历}

\begin{tabularx}{\linewidth}{ @{}l r@{} }
\textbf{RoboMaster 竞赛：串联弹性腿平衡步兵机器人控制系统} & \hfill \href{https://Particlela.github.io/project/}{2024年10月 - 至今} \\[3.75pt]
\multicolumn{2}{@{}X@{}}{基于 STM32H7 平台为 RoboMaster 串联弹性腿轮腿平衡步兵机器人开发全栈嵌入式控制系统，基于五杆腿机构、面向高动态运动需求。核心控制层采用\textbf{可变增益 LQR 平衡控制策略}，设计 10 状态 LQR 控制器，将增益矩阵与前馈项表示为左右虚拟腿长度的二元三次多项式函数，进行在线实时插值，保证在不同腿构型下稳定平衡。决策与行为调度层运行\textbf{基于优先级的分层状态机}，覆盖 Dead、Recovery、Mature、Flight、Jump、Onestep、Twostep、Creep 等多种工作状态，支持越障跳跃与单/双步爬升等复杂动作的有序切换与执行。能量层面集成超级电容与缓冲电容形成两级能量缓冲，并引入\textbf{基于二次模型的功率限幅器}，根据裁判系统实时功率限制与剩余能量状态动态缩放量速度与偏航角速度指令，在有限功率预算下平衡机动性与系统安全性，最终形成平衡控制、行为管理与能量调度一体化的机器人控制闭环。}  \\
\end{tabularx}

\begin{tabularx}{\linewidth}{ @{}l r@{} }
\textbf{RoboMaster 竞赛：串联弹性腿哨兵机器人控制系统} & \hfill \href{https://Particlela.github.io/project/}{2025年8月 - 至今} \\[3.75pt]
\multicolumn{2}{@{}X@{}}{在步兵平衡机器人框架基础上升级为面向自主导航的哨兵机器人，沿用同一套串联弹性腿硬件平台与底层控制架构，将控制重心由人工操作转向自主决策与高层战场机动。为实现无人化运行，系统新增\textbf{底盘导航指令接口}，接收上层 NUC 的运动参考指令，使哨兵完全摆脱遥控或键鼠输入依赖、独立执行导航任务。哨兵机器人不仅继承了步兵的高动态平衡能力，还具备在未知环境中自主规划与执行复杂机动的完整智能化能力。}  \\
\end{tabularx}

\begin{tabularx}{\linewidth}{ @{}l r@{} }
\textbf{双向四开关 Buck-Boost 缓冲电容器控制系统（FSBB）} & \hfill \href{https://Particlela.github.io/project/}{2026年6月 - 至今} \\[3.75pt]
\multicolumn{2}{@{}X@{}}{基于 STM32G4 MCU，设计并实现面向 RoboMaster 移动机器人的高带宽缓冲电容能量管理控制器。主功率级采用双向四开关 Buck-Boost 变换器拓扑，可动态吸收与释放能量，有效抑制高动态机动过程中的瞬时母线电压跌落。控制层运行\textbf{基于状态机的多环路控制策略}，集成电压比前馈补偿与三个独立 PI 环路（恒功率充电、电容电压调节、母线电压支撑），使控制器能根据母线状态在充电、调压与放电模式间平滑切换；同时使用多通道 ADC 与 DMA 高速采集输入电压、电容电压、输入电流、输出电流与电容电流，并经精度标定与互补滤波为各闭环提供高精度、低延迟的实时状态反馈。安全机制集成功率缓启、过压/欠压/过流阈值保护，以及 CAN 通信超时或异常工况下的自动关停与恢复逻辑，确保在复杂电气环境下硬件可靠运行。关键性能指标：电容电压工作范围 11 V–28 V，最大充电功率 190 W，最大输出/补充功率 1200 W，电容电流限值 70 A，输入电压范围 15 V–28 V。}  \\
\end{tabularx}

\begin{tabularx}{\linewidth}{ @{}l r@{} }
\textbf{HW-Components：通用机器人控制库} & \hfill \href{https://Particlela.github.io/project/}{2024年10月 - 至今} \\[3.75pt]
\multicolumn{2}{@{}X@{}}{作为核心贡献者，架构并维护了面向 RoboMaster 及通用移动机器人平台的可复用机器人控制库 HW-Components 的嵌入式通信层。}  \\
\end{tabularx}

\begin{tabularx}{\linewidth}{@{}X@{}}
\begin{itemize}[nosep,after=\strut, leftmargin=1em, itemsep=3pt,label=--]
\item \textbf{多协议测距传感器通信}：实现 STM32 与多种测距模块（超声波、ToF、红外等）的单向与双向数据交互，为不同机器人构型提供灵活的感知接口。
\item \textbf{裁判系统串口链路}：维护并优化与 RoboMaster 裁判系统基于 UART 的通信协议，确保机器人状态、受击点、开火权限与功率限制等比赛关键数据的可靠接收。
\item \textbf{基于 CAN 的能量管理}：开发超级电容与缓冲电容子系统的双向 CAN 通信，实现实时功率监测、能量缓冲控制与动态功率分配。
\item \textbf{杠杆臂补偿}：实现 IMU 安装位置与机器人旋转中心物理偏移所引入的运动学误差补偿，提升快速姿态变化过程中的状态估计精度。
\end{itemize}
\end{tabularx}

%----------------------------------------------------------------------------------------
%	PUBLICATIONS
%----------------------------------------------------------------------------------------
% \section{Publications}
% \begin{refsection}[citations.bib]
% \nocite{*}
% \printbibliography[heading=none]
% \end{refsection}

%----------------------------------------------------------------------------------------
%	SKILLS
%----------------------------------------------------------------------------------------
\section{专业技能}
\begin{tabularx}{\linewidth}{@{}l X@{}}
编程语言 & \normalsize{C/C++, Python, MATLAB}\\
嵌入式平台 & \normalsize{STM32 (H7/G4), NUC, Raspberry Pi}\\
控制理论 & \normalsize{PID, LQR, MPC, EKF, 状态机, 自适应控制}\\
通信协议 & \normalsize{CAN, UART, SPI, I2C, DMA}\\
开发工具 & \normalsize{CMake, Git, Linux, ROS2, OpenCV, LaTeX}\\
\end{tabularx}

% \vfill
% \center{\footnotesize Last updated: \today}

\end{document}